\section{Final checklist}

This chapter introduced empirical distributions and descriptive statistics as
ways of summarizing observed data. The most important idea is the distinction
between theoretical objects from probability models and empirical objects
computed from finite samples.

\subsection*{Conceptual understanding}

You should be able to explain the following distinctions:
\begin{itemize}
    \item a probability law versus a finite dataset;
    \item a random variable \(X_i\) before observation versus an observed value
    \(x_i\) after observation;
    \item theoretical probability \(P(X\in A)\) versus empirical probability
    \(P_n(A)\);
    \item theoretical CDF \(F_X\) versus empirical CDF \(F_n\);
    \item density versus histogram;
    \item theoretical expectation \(\E[X]\) versus sample mean \(\bar{x}\);
    \item theoretical variance \(\Var(X)\) versus empirical or corrected sample
    variance;
    \item theoretical quantiles versus empirical quantiles.
\end{itemize}

\subsection*{Definitions and formulas}

You should know the meaning of the following formulas:
\[
P_n(A)=\frac1n\sum_{i=1}^n\mathbf{1}_{\{x_i\in A\}},
\]
\[
F_n(x)=\frac1n\sum_{i=1}^n\mathbf{1}_{\{x_i\leq x\}},
\]
\[
\bar{x}=\frac1n\sum_{i=1}^n x_i,
\]
\[
v_n=\frac1n\sum_{i=1}^n(x_i-\bar{x})^2,
\]
\[
s^2=\frac1{n-1}\sum_{i=1}^n(x_i-\bar{x})^2,
\]
\[
IQR=Q_3-Q_1.
\]
You should also understand why the denominator \(n\) describes the empirical
distribution, while the denominator \(n-1\) is connected with later inferential
use.

\subsection*{Typical tasks}

After studying this chapter, you should be able to:
\begin{itemize}
    \item build a frequency table from discrete data;
    \item compute relative frequencies and interpret them as empirical
    probabilities;
    \item compute values of the empirical CDF by hand for a small dataset;
    \item explain why the empirical CDF is a step function;
    \item draw or interpret a histogram;
    \item explain why bin choice affects a histogram;
    \item compute the sample mean, empirical variance, standard deviation,
    median, quartiles, and IQR;
    \item identify how outliers affect means and standard deviations;
    \item explain why computer-generated samples help illustrate the relationship
    between a model and data;
    \item compare empirical summaries with theoretical quantities without
    confusing them.
\end{itemize}

\subsection*{Common mistakes}

Avoid the following mistakes:
\begin{itemize}
    \item saying that a sample is itself a normal distribution, exponential
    distribution, or other theoretical law;
    \item treating \(\bar{x}\) as if it were automatically equal to \(\E[X]\);
    \item treating an empirical histogram as if it were exactly a density;
    \item forgetting that histograms depend on bin choices;
    \item ignoring the difference between \(n\) and \(n-1\) in variance formulas;
    \item removing outliers mechanically without understanding their context;
    \item treating simulation as proof rather than illustration.
\end{itemize}

\subsection*{Questions for self-explanation}

A good test of understanding is whether you can answer these questions clearly:
\begin{enumerate}
    \item What is the empirical distribution of a dataset?
    \item Why is every finite empirical distribution discrete?
    \item What does \(F_n(x)\) measure?
    \item How is a histogram related to a density, and how is it different?
    \item Why can two samples from the same model have different sample means?
    \item What is the difference between describing a sample and estimating a
    population quantity?
    \item Why do computer-generated samples help us understand probability models?
    \item What will change in the next chapter when statistics such as \(\bar{X}\)
    are treated as random variables?
\end{enumerate}

\begin{summarybox}
Empirical distributions and descriptive statistics translate the language of
probability into the language of observed data. They show how the concepts of
distribution, center, spread, and quantiles begin to work when only a finite
sample is available. The next chapter studies the randomness of these summaries
under repeated sampling.
\end{summarybox}
